
\textbf{\SGAInlineRefTarget{sga3:xxiv:proposition:3-13}Proposition 3.13.} \emph{Let \(S\) be a scheme, \(G\) a semisimple adjoint (resp. simply connected) \(S\)-group, \(B \supset T\) a
Killing couple of \(G\), \(\operatorname{Dyn}(G)\) the Dynkin \(S\)-scheme}

\emph{of \(G\). There exists a canonical isomorphism of \(S\)-group schemes}

\[
  T\xrightarrow{\sim}
  \prod_{\operatorname{Dyn}(G)/S}
  \mathbb G_{m,\operatorname{Dyn}(G)}.
\]

\emph{(Recall (cf.~Exp. II, § 1) that the right-hand side is by definition the \(S\)-functor which to \(S' \to S\) associates
\(G_{m}(\operatorname{Dyn}(G) \times_{S} S')\), or, what amounts to the same, \(G_{m}(Dyn(G_{S'}))\).)}

\emph{First proof.} Let us do it for simplicity in the adjoint case. Consider the composite morphism

\[
  T\longrightarrow\prod_{D/S}T_D
  \longrightarrow\prod_{D/S}\mathbb G_{m,D}.
\]

where the first morphism is the canonical morphism, the second is \(\prod_{D/S} \alpha_{D}\) (we have written \(D =
\operatorname{Dyn}(G)\)). To verify that this morphism is an isomorphism, one may assume \(G\) split; now, in this case, this is none
other than the morphism \(T \to (G_{m, S})^{\Delta}\) with components \(\alpha\), for \(\alpha \in \Delta\), and the latter is
an isomorphism (\SGARefLink{sga3:xxii:rem:4.3.8}{Exp. XXII, 4.3.8}).

\emph{Second proof.} By \SGARefLink{sga3:xxiv:2.8}{2.8}, \SGARefLink{sga3:xxiv:3.5}{3.5} and \SGARefLink{sga3:xxiv:3.11}{3.11}, one may assume that \(G = Q-\acute{E}p_{S'/S}(R)\), \(T\) being the canonical maximal
torus. On \(S'\), one has by \SGARefLink{sga3:xxii:rem:4.3.8}{Exp. XXII 4.3.8}, an isomorphism \(T_{S'} \xrightarrow{\sim} (G_{m, S'})^{\Delta}\), defined by
the simple roots (resp. the simple coroots). The group \(E\) operates on the right-hand side by permutation of \(\Delta\).
Now the bundle associated to \(S'/S\) by \(E \to \operatorname{Aut}(\Delta)\) is \(\operatorname{Dyn}(G)\) (\SGARefLink{sga3:xxiv:3.7}{3.7}), and one concludes at once.

Using Prop. \SGARefLink{sga3:xxiv:proposition:8-4}{8.4} of the appendix, one deduces:

\textbf{\SGAInlineRefTarget{sga3:xxiv:corollary:3-14}Corollary 3.14.} \emph{Under the preceding conditions, one has}

\[
  H^1(S,T)\xrightarrow{\sim}
  H^1\bigl(\operatorname{Dyn}(G),\mathbb G_m\bigr)
  \xrightarrow{\sim}\operatorname{Pic}\bigl(\operatorname{Dyn}(G)\bigr).
\]

\emph{In particular, \(H^{1}(S, T) = 0\) when \(S\) is semi-local.}

\textbf{\SGAInlineRefTarget{sga3:xxiv:remark:3-15}Remark 3.15.} Let \(S\) be a scheme, \(G\) (resp. \(G'\)) an \(S\)-reductive group, \(B \supset T\) (resp. \(B' \supset T'\)) a
Killing couple of \(G\) (resp. \(G'\)), \(u \in Isomext(G, G')(S)\). Set (cf.~\SGARefLink{sga3:xxiv:2.10}{2.10})

\[
  P=\underline{\operatorname{Isomint}}_u(G,B,T;G',B',T').
\]

this is a principal homogeneous bundle under \(T_{ad}\) (by \((f, t) \mapsto f \circ int(t)\)).

Let on the other hand \(D = \operatorname{Dyn}(G) = \operatorname{Dyn}(G')\) (identified via \(u\) (\SGARefLink{sga3:xxiv:3.5}{3.5})), and let \(L = Isom_{D}(g_{D}, g'_{D})\) be the
principal homogeneous bundle under \(G_{m,D}\) defined by the invertible \(\mathcal O_D\)-module

\[
  \operatorname{Hom}_{\mathcal O_D}(\mathfrak g_D,\mathfrak g'_D)
  =\mathfrak g'_D\otimes(\mathfrak g_D)^\vee .
\]

Each \(f \in P(S')\) defines an isomorphism of \(g_{D} \otimes_{O_{S}} O_{S'}\) onto \(g'_{D} \otimes_{O_{S}} O_{S'}\), whence
a canonical morphism

\[ P \to \prod_{D/S} L. \]

This morphism is an isomorphism, compatible with the isomorphism of operator groups

\[
  T_{\mathrm{ad}}\xrightarrow{\sim}
  \prod_{D/S}\mathbb G_{m,D}.
\]

defined above. Indeed, it suffices to verify it in the case where the two groups are pinned, where it is easy.

It follows in particular that in the isomorphism

\[ H^{1}(S, T_{ad}) \xrightarrow{\sim} \operatorname{Pic}(\operatorname{Dyn}(G)) \]

of \SGARefLink{sga3:xxiv:corollary:3-14}{3.14}, the class of the bundle \(P\) is transformed into \(c\ell(g'^{D}) - c\ell(g_{D})\). The bundle \(P\) is therefore
trivial if and only if \(g'^{D}\) and \(g_{D}\) are isomorphic.

If \((G, B, T)\) is quasi-splittable, for example if one takes for \(G\) the group \(G'_{q-\acute{e}p.}\), with its canonical
Killing couple, it follows that the image of the class of \(P\) is none other than the obstruction to the quasi-splitting
of \(G'\) defined in \SGARefLink{sga3:xxiv:3.9}{3.9}.

\subsection*{\SGAInlineRefTarget{sga3:xxiv:subsection:3-16}3.16. Symmetries}

\SGAInlineRefTarget{sga3:xxiv:subsubsection:3-16-1}\subsubsection*{3.16.1.}

Let \(S\) be a scheme, \(G\) an \(S\)-reductive group, \(B \supset T\) a Killing couple of \(G\). Write \(D = \operatorname{Dyn}(G)\). Recall (\SGARefLink{sga3:xxii:prop:5.9.2}{Exp.
XXII, 5.9.2}) that there exists a unique Borel subgroup \(B^{-}\) of \(G\) such that \(B \cap B^{-} = T\). If \(X \in \Gamma(D,
g_{D})^{\times}\) defines a quasi-pinning of \(G\) relative to \((B, T)\) (\SGARefLink{sga3:xxiv:3.9}{3.9}), then \(Y = -X^{-1} \in \Gamma(D,
(g_{D})^{\vee})^{\times}\) defines a quasi-pinning of \(G\) relative to \((B^{-}, T)\); one says that this is the \emph{opposite
quasi-pinning}.

If \(R\) is a pinned reduced root datum and if \(E\) is an \(R\)-pinning of the \(S\)-reductive group \(G\), one defines an
\(R\)-pinning \(E^{-}\) said to be \emph{opposite to \(E\)} in the following way: one keeps the same maximal torus \(T\), one takes
the opposite of the isomorphism \(D_{S}(M) \xrightarrow{\sim} T\), and one ``pins'' by \(Y_{\alpha} = -X^{-1}_{\alpha} \in
\Gamma(S, g^{-\alpha})^{\times}\), for \(\alpha \in \Delta\). The quasi-pinning underlying \(E^{-}\) is the quasi-pinning
opposite to the quasi-pinning underlying \(E\).

\textbf{Remark.} In the notation of \SGARefLink{sga3:xx:3.1}{Exp. XX, 3.1}, if one sets

\[
  w_\alpha(X_\alpha)
  =\exp(X_\alpha)\exp(-X_\alpha^{-1})\exp(X_\alpha).
\]

one has \(w_{\alpha}(X_{\alpha}) = w_{-\alpha}(Y_{\alpha})\) (loc. cit., 3.1 (vi)).

\textbf{\SGAInlineRefTarget{sga3:xxiv:proposition:3-16-2}Proposition 3.16.2.} \emph{Let \(S\) be a scheme, \(G\) an \(S\)-reductive group.}

\emph{\SGAInlineRefTarget{sga3:xxiv:proposition:3-16-2:item:i}(i) Let \(T\) be a maximal torus of \(G\); there exists a unique}

\[
  i_T\in
  \left(
    \underline{\operatorname{Aut}}_{S\text{-}\mathrm{gr.}}(G,T)/T_{\mathrm{ad}}
  \right)(S)
  \subset
  \underline{\operatorname{Aut}}_{S\text{-}\mathrm{gr.}}(T).
\]

\emph{such that \(i_{T}(t) = t^{-1}\) for every section \(t\) of \(T\).}

\emph{\SGAInlineRefTarget{sga3:xxiv:proposition:3-16-2:item:ii}(ii) Let \((B, T)\) be a Killing couple of \(G\); there exists a unique section}

\[
  w_{B,T}\in
  \bigl(\operatorname{Norm}_G(T)/T\bigr)(S)=W_G(T)(S).
\]

\emph{such that \(int(w_{B, T})(B) = B^{-}\) (with the obvious abuse of language).}

\emph{\SGAInlineRefTarget{sga3:xxiv:proposition:3-16-2:item:iii}(iii) Let \(Q = (B, T, X)\) be a quasi-pinning of \(G\), \(Q^{-} = (B^{-}, T, Y)\) the opposite quasi-pinning; there exists a
unique inner automorphism of \(G\)}

\[ n_{Q} \in Norm_{ad(G)}(T_{ad})(S) \subset \underline{\operatorname{Aut}}_{S\text{-}\mathrm{gr.}}(G) \]

\emph{such that \(n_{Q}(Q) = Q^{-}\), that is, \(n_{Q}(T) = T\), \(n_{Q}(B) = B^{-}\), \(n_{Q}(X) = Y\).}

\emph{\SGAInlineRefTarget{sga3:xxiv:proposition:3-16-2:item:iv}(iv) Let \((R, E)\) be a pinning of \(G\), \((R, E^{-})\) the opposite pinning; there exists a unique automorphism of \(G\)}

\[
  u_E\in
  \underline{\operatorname{Aut}}_{S\text{-}\mathrm{gr.}}(G,T)
  \subset
  \underline{\operatorname{Aut}}_{S\text{-}\mathrm{gr.}}(G).
\]

\emph{such that \(u_{E}(E) = E^{-}\), i.e.~\(u_{E}(t) = t^{-1}\) for every section \(t\) of \(T\), and \(Ad(u_{E}) X_{\alpha} =
Y_{\alpha}\) for every \(\alpha \in \Delta\).}

\emph{Proof.} \SGARefLink{sga3:xxiv:proposition:3-16-2:item:ii}{(ii)} follows from \SGARefLink{sga3:xxii:corollary:5-5-5:item:ii}{Exp. XXII, 5.5.5 (ii)}, then \SGARefLink{sga3:xxiv:proposition:3-16-2:item:iii}{(iii)} follows from \SGARefLink{sga3:xxiv:3.10}{3.10}, and \SGARefLink{sga3:xxiv:proposition:3-16-2:item:iv}{(iv)} follows from \SGARefLink{sga3:xxiii:theorem:4-1}{Exp. XXIII, 4.1}.
Finally to prove \SGARefLink{sga3:xxiv:proposition:3-16-2:item:i}{(i)}, one may assume \(G\) pinned. Existence follows from \SGARefLink{sga3:xxiv:proposition:3-16-2:item:iv}{(iv)} for example, uniqueness from the fact that
an automorphism of \(G\) inducing the identity on \(T\) is given by a section of \(T_{ad}\) (\SGARefLink{sga3:xxiv:2.11}{2.11}).

\textbf{\SGAInlineRefTarget{sga3:xxiv:corollary:3-16-3}Corollary 3.16.3.} \emph{One has}

\[
  i_T^2=w_{B,T}^2=n_Q^2=u_E^2=e.
\]

\emph{Moreover, \(i_{T}\) (resp. \(u_{E}\)) is \(\neq e\) if \(G \neq e\), and \(w_{B, T}\) (resp. \(n_{Q}\)) is \(\neq e\) if \(G\) is not a
torus.}

\textbf{\SGAInlineRefTarget{sga3:xxiv:corollary:3-16-4}Corollary 3.16.4.} \emph{In the situation of \SGARefLink{sga3:xxiv:proposition:3-16-2:item:iii}{(iii)} (resp. \SGARefLink{sga3:xxiv:proposition:3-16-2:item:iv}{(iv)}), \(n_{Q}\) projects onto \(w_{B, T}\) (resp. \(u_{E}\) projects
onto \(i_{T}\)) by the canonical morphism}

\[ Norm_{ad(G)}(T_{ad}) \to W_{ad(G)}(T_{ad}) \simeq W_{G}(T), \]

\emph{resp.}

\[
  \underline{\operatorname{Aut}}_{S\text{-}\mathrm{gr.}}(G,T)
  \longrightarrow
  \underline{\operatorname{Aut}}_{S\text{-}\mathrm{gr.}}(G,T)/T_{\mathrm{ad}}.
\]

\textbf{\SGAInlineRefTarget{sga3:xxiv:corollary:3-16-5}Corollary 3.16.5.} \emph{The preceding definitions are compatible with base change, and are functorial under isomorphism
(in an obvious sense).}

\textbf{\SGAInlineRefTarget{sga3:xxiv:proposition:3-16-6}Proposition 3.16.6.} \emph{\SGAInlineRefTarget{sga3:xxiv:proposition:3-16-6:item:i}(i) One can define uniquely for each reductive group \(G\) over a scheme \(S\) an element}

\[ s_{G} \in Autext(G)(S) \]

\emph{in such a way that this construction is functorial in \(G\) under isomorphism, is compatible with base change, and is
such that whenever \(T\) is a maximal torus of the \(S\)-reductive group \(G\), \(s_{G}\) is the image of \(i_{T}\) under the
canonical morphism}

\[
  \underline{\operatorname{Aut}}_{S\text{-}\mathrm{gr.}}(G,T)/T_{\mathrm{ad}}
  \longrightarrow\underline{\underline{\operatorname{Autext}}}(G).
\]

\emph{\SGAInlineRefTarget{sga3:xxiv:proposition:3-16-6:item:ii}(ii) One has \(s^{2}_{G} = e\), and \(s_{G}\) is a central element of \(Autext(G)\).}

\emph{\SGAInlineRefTarget{sga3:xxiv:proposition:3-16-6:item:iii}(iii) Under the conditions of \SGARefLink{sga3:xxiv:proposition:3-16-2:item:ii}{3.16.2 (ii)}, if one identifies \(\underline{\operatorname{Aut}}_{S\text{-}\mathrm{gr.}}(G, B, T) / T_{ad}\) with
\(Autext(G)\) (\SGARefLink{sga3:xxiv:2.2}{2.2}), one has}

\[
  w_{B,T}i_T=i_Tw_{B,T}=s_G.
\]

\emph{\SGAInlineRefTarget{sga3:xxiv:proposition:3-16-6:item:iv}(iv) Under the conditions of \SGARefLink{sga3:xxiv:proposition:3-16-2:item:iv}{3.16.2 (iv)}, if one identifies \(\operatorname{Aut}_{S-gr.-\acute{e}p}(G, R, E)\) with
\(Autext(G)\) (\SGARefLink{sga3:xxiv:1.3:item:iii}{1.3 (iii)}), one has}

\[
  n_Eu_E=u_En_E=s_G.
\]

\emph{Proof.} \SGARefLink{sga3:xxiv:proposition:3-16-6:item:i}{(i)} is proved without difficulty by descent. On the other hand, since \(i_{T}\) is evidently a central section of
square \(e\) in \(\underline{\operatorname{Aut}}_{S\text{-}\mathrm{gr.}}(T)\), \SGARefLink{sga3:xxiv:proposition:3-16-6:item:ii}{(ii)} follows immediately; \SGARefLink{sga3:xxiv:proposition:3-16-6:item:iii}{(iii)} is a consequence of \SGARefLink{sga3:xxiv:proposition:3-16-6:item:iv}{(iv)} by descent.
Finally, under the conditions of \SGARefLink{sga3:xxiv:proposition:3-16-6:item:iv}{(iv)}, it is clear that \(n_{E} u_{E} = u_{E} n_{E}\) and that this automorphism of \(G\)
respects the pinning; modulo the identification made, it is therefore equal to its image in \(Autext(G)\); but \(n_{E}\) is
inner and \(u_{E}\) projects onto \(s_{G}\).

\textbf{\SGAInlineRefTarget{sga3:xxiv:remark:3-16-7}Remark 3.16.7.} \SGARefLink{sga3:xxiv:proposition:3-16-6:item:i}{(i)} One determines \(s_{G}\) explicitly in each case of the classification thanks to \SGARefLink{sga3:xxiv:proposition:3-16-6:item:iii}{(iii)}: for each
irreducible pinned root datum it suffices to compose the symmetry through the origin with the symmetry in the Weyl group
(i.e.~the element \(w_{0}\) of the Weyl group such that \(w_{0}(\Delta) = -\Delta\)). One finds the following results: one
has \(s_{G} = e\) except for \(A_{n}\) (\(n \geqslant 2\)), \(D_{n}\) (\(n\) odd) and \(E_6\), in which case \(s_{G}\) is the unique
non-trivial ``outer automorphism''.

\begin{enumerate}
\def\labelenumi{(\roman{enumi})}
\setcounter{enumi}{1}
\item
  The automorphism \(u_{E}\) is the one used to construct ``the compact real forms'' in the theory of semisimple Lie
  algebras.
\end{enumerate}

\textbf{\SGAInlineRefTarget{sga3:xxiv:remark:3-16-8}Remark 3.16.8.} We defined in \SGARefLink{sga3:xxiv:subsubsection:3-16-1}{3.16.1} an involution in the \(S\)-scheme
\[
Q = Isomint(G_{q-\acute{e}p.}, G)
\]
of quasi-pinnings of \(G\) (cf.~\SGARefLink{sga3:xxiv:3.12.1}{3.12.1}); by transport of structure from \(G_{q-\acute{e}p.}\) to \(G\), one sees at once that
this involution is given by the action of an element of \(ad(G_{q-\acute{e}p.})(S)\): the element \(n_{0}\) defined (\SGARefLink{sga3:xxiv:proposition:3-16-2:item:iii}{3.16.2
(iii)}) by the canonical quasi-pinning of \(G_{q-\acute{e}p.}\).

In the same way, one has defined an involution in the \(S\)-scheme

\[ P = Isom_{S-gr.}(\acute{E}p_{S}(R), G) \]

of \(R\)-pinnings of \(G\) (cf.~\SGARefLink{sga3:xxiv:1.20}{1.20}). Arguing as before, one sees that this involution is given by the action of the
automorphism \(u_{0}\) of \(\acute{E}p_{S}(R)\) defined (\SGARefLink{sga3:xxiv:proposition:3-16-2:item:iv}{3.16.2 (iv)}) by the canonical pinning of \(\acute{E}p_{S}(R)\).
